% !TEX program = lualatex
\documentclass[paper=a4, fontsize=10.5pt]{jlreq}

% 数式・記号
\usepackage{amsmath,amssymb}
\usepackage{bm}

% 画像・図表
\usepackage{graphicx} % [dvipdfmx] は不要
\usepackage{booktabs}
\usepackage{dcolumn}
\usepackage{lscape}
\usepackage[flushleft]{threeparttable}

% レイアウト・書式
\usepackage{times}
\usepackage[top=30truemm,bottom=30truemm,left=25truemm,right=25truemm]{geometry}
\usepackage{authblk}
\usepackage{comment}
\usepackage{url}
\usepackage{natbib}


\title{Results of Baseline and Bi-weekly Surveys }
\author{Kazushi Takahashi, Wataru Kodama, Yukichi Mano}
\date{
	\today
}

\begin{document}
	\maketitle
	%\newpage
\section{Descriptive statistics}  
天水地域

各郡のディレクターに全ての稲作コミュニティをリストしてもらう（Achiase 11; Bebiani-Anhwiaso-Bekwai 38; Biakoye 19; Lower Manya Krobo 4; Pru West 11; Sefwi Wiawso 27; Sene West 14; Krachi West 20）。
全ての稲作コミュニティについて、担当AEAに各コミュニティと連絡をとってもらい（幹部などに偏らないよう、なるべくランダムに）8名の稲作農家をリストしてもらう。
各AEAにつき２つの稲作コミュニティをランダムに抽出する。各サンプルコミュニティの農家リストからランダムに5名をサンプル農家を抽出し、残り3名をリザーブ農家とする。ただし、担当の稲作コミュニティがひとつしかないAEAについては8名全てをサンプル農家として抽出。


灌漑地域

各郡のディレクターからGRIPを通じてKpongとWetaの全稲作農家のリストを入手。
Kpong（約2800農家）は15のbranch canalのうちの１つのSLLCA（約390農家）をランダムに選び、農家リストからランダムに10名をサンプル農家、5名をリザーブとして抽出。(当初、別の2つのbranch canal、SLLCとNLLCAも抽出し、それぞれ6名と8名の農家を選んでいたが、Dr. Adakuから各農家の村があまりに分散しているとのことで、SLLCAに絞ることになった)
Weta（約1200農家）は農家リストからランダムに10名をサンプル農家、10名をリザーブとして抽出。

\begin{table}[hbtp]
	\centering
	\begin{threeparttable}
		\caption{Number of Observations by District}
		
		% ここで列の定義を自分で行う
		\begin{tabular}{llccc}
			\hline
			\input{tmp/obs_data.tex} % Stataが出力した中身だけを読み込む
			
		\end{tabular}
		
	\end{threeparttable}
\end{table}

	\begin{table}[hbtp]
		\centering
		\scriptsize
		\begin{threeparttable}
			\caption{Demographic Characteristics}
			
			\input{tmp/f_balance1.tex}
			
		\end{threeparttable}
	\end{table}

	\begin{table}[hbtp]
	\centering
	\scriptsize
	\begin{threeparttable}
		\caption{Farm Characteristics}
		
		\input{tmp/f_balance2.tex}
		
	\end{threeparttable}
\end{table}

 \begin{landscape} 
\begin{table}[hbtp]
	\centering
	\scriptsize
	\begin{threeparttable}
		\caption{Overall characteristics of AEAs}
		
		\input{tmp/a_basic_balance.tex}
		
	\end{threeparttable}
\end{table}
\end{landscape}

 \begin{landscape} 
	\begin{table}[hbtp]
		\centering
		\scriptsize
		\begin{threeparttable}
			\caption{Prosocialness}
			
			\input{tmp/a_int_balance.tex}
			
		\end{threeparttable}
	\end{table}

\end{landscape}

\begin{landscape} 
\begin{figure}
	
	\hspace{3cm}
	\includegraphics[scale=1,clip]{tmp/prosocial.eps}
	\caption{Prosocialness}
\end{figure}
\end{landscape}


 \begin{landscape} 
	\begin{table}[hbtp]
		\centering
		\scriptsize
		\begin{threeparttable}
			\caption{Locus of control}
			
			\input{tmp/a_loc_balance.tex}
			
		\end{threeparttable}
	\end{table}
\end{landscape}

\begin{landscape} 
	\begin{figure}
		
		\hspace{3cm}
		\includegraphics[scale=1,clip]{tmp/locus1.eps}
		\caption{Locus of control (agriculture)}
	\end{figure}
\end{landscape}

\begin{landscape} 
	\begin{figure}
		
		\hspace{3cm}
		\includegraphics[scale=1,clip]{tmp/locus2.eps}
		\caption{Locus of control (general)}
	\end{figure}
\end{landscape}

	
\end{document}
